\magnification=\magstephalf
\input notesmac.tex
\nopagenumbers

\centerline{\titlefont CZ3106 Symbolic Computing, Lab 1}
\bigskip
\line{Dr.~Wang Jian-Sheng \hfill For the week 13-18 Jan 03, due one week after the lab}
\bigskip
%----------------------------------------------------------------------- 

\problem Try to familiarize with the notebook interface.  Read the 
first few sections of the Mathematica Book.  Use the help system to
find out the command {\tt N[ ]}.  Type some text (as document) and
enter a command.  Change the font/color of the text.  Change the 
stylesheet.  Try cut and paste, etc.  You need not hand-in anything  
for this problem. 

\problem We want to be familiar ourselves to arithmetics in  
Mathematica.  We can use it as a pocket calculator.  
Consider entering the following mathematical expressions
to the system and watch the reaction of the system:
$$\halign{\h$#$\hfil&\h$#$\hfil&\h$#$\hfil\cr
1 + 2, &  2-3\times5 -1/2, & 2/3 - 0.5, \cr
\sqrt{4}, & i^2, &  2^{10}.\cr}$$

\problem All symbolic systems 
can handle arbitrarily large integers without overflow. Try the
following:
$$\halign{\h$#$\hfil&\h$#$\hfil&\h$#$\hfil\cr
2^{100}, & 30!, & 123456789*987654321. \cr}$$

\problem Mathematica can do arbitrarily precise floating-point
computation. 
{\tt N[ ]} convert an exact number into accurate precision floating-point
value.  Find the numerical values of the following numbers accurate
to 30 decimal places (quad precision):
$$\halign{\h$#$\hfil&\h$#$\hfil&\h$#$\hfil\cr
\pi, &  e, & \sqrt{2}, \cr
355/113, & \sin {\pi\over7}, & e^{\pi \sqrt{163}}.  \cr}$$
In Mathematica, make the assignment {\tt pi = N[Pi,30]}.  What is the
digits of precision of {\tt 2*pi}, and what is {\tt 2.0*pi}?  Why the
second expression gives you only six digits?  How to remove the symbol
{\tt pi} from the system when you have finished using it?

\problem In this exercise we like to see how Mathematica deal with
symbols.  What happens when the following mathematical expressions are
typed in into symbolic systems:
$$\halign{\h$#$\hfil&\h$#$\hfil&\h$#$\hfil\cr
\pi, &  x+y, & x-x, \cr
2x\times x+1 & (1+x)^2 & \sin {\pi\over 2}.  \cr}$$
Note that $\pi$ should be typed in as the constant {\tt Pi}, not a numerical
approximation.          

\problem Consider the cubic equation
$$ f(x) = x^3 + 6 x^2 + 9 x + 1 = 0. $$
\itemitem{(a)}  Make a plot to see how many real roots this equation
has.  Read from the plot the rough values of the solutions.
\itemitem{(b)} Solve the equation numerically.
\itemitem{(c)} Solve the equation symbolically.  The resulting expressions
are complicated with imaginary number $i$.   But we know the solution 
has to be real (from the plot).  Simplify the expressions if you can. 


\problem Plot
the curve (using {\tt ParametricPlot[ ]}) 
given by $r = \theta$ (polar coordinates).  Plot
two curves given by $r = \theta$ and $r = 1$ on the same plot, using
two different color. 
 
%% 
%% Map[ComplexExpand, Solve[x^3+6x^2+9x+1==0, x], {3}]
%% 
%%

\bye
